\section{Quantitative Analysis}
In this section, I calibrate the quantitative model using administrative data from France. I then use it to quantitatively test my model against the empirical evidence on wage passthrough and layoffs across firms and tenure.
\subsection{Solving the model}
%Core idea: current model still untractable. I argue, based on the theoretical propositions, that the contracts converge. I also refer to the data oon wage growth across tenure.
The tenure-based formulation of the problem allows me to work with a discrete, although expanding, state-space. Finally, to make the problem fully tractable, I note that wages in contracts tend to converge. The result uses the propositions derived in the previous subsection.
One corollary of Proposition~\ref{prop:targetwage} is that wages between the cohorts of the same quality will converge since they have the same target wage. We can also expect quality across cohorts to converge:
\begin{itemize}
    \item Lower quality cohorts (compared to other cohorts) will have lower target wages
    \item Over time, low quality cohorts become low value cohorts
    \item Following Proposition~\ref{prop:qualityconv}, lower value cohorts will be the first to get laid off
    \item Thus, lower quality cohorts will over time catch up in quality and, thus, also in wages
\end{itemize}
Given that the contracts should converge, the question remains how quickly it happens, or, from a practical perspective, how many cohorts should we keep track of? Empirically, I show in Appendix~\ref{wagegrowthK} that wage growth in France stagnates after about 10 years of tenure.  I use this to justify restricting attention to finite and constant $K\leq 10$ for all the firms in my quantitative exploration. Note that this is only an approximation since, as $K\rightarrow\infty$, the model approaches the problem described in Lemma~\ref{firmproblem}. \\
Another complication, common for models of dynamic contracts, is large action space due to $v'_{y',k}$ expanding alongside the number of possible productivity states. In Appendix~\ref{dual} I show that the problem can be solved in its dual form by choosing future marginal utilities, which are constant across productivity realizations, instead of promised values.
\\
Lastly, I solve for the block recursive equilibrium of this model. The equilibrium exists if all the workers, when hired, are hired on the exact same value $v_0$ ( see Appendix~\ref{blocrecur} for the details). The rest of the value that the workers are owed from searching in higher submarkets $v$ is distributed to them immediately upon hiring via a sign-on wage. I set the initial value $v_0$ to the unemployment value $U$.
\subsection{Model specification}
%I quantify the model by matching moments of the matched employer-employee data from France between 2009-2019. \\
%Explain how I choose my parameters 
%I start by discussing the parameters I set externally. \\
I choose to work with the yearly frequency as contracts in France are rarely updated more than once a year, providing me with a clear way to track cohorts. I set the discounting rate of $\beta = 0.96$ to match an annual interest rate of $4\%$. \\ 
I work with the logarithmic utility as I need both the inverse derivative $u'(u^{-1}(v))$ function to be strictly increasing, which, within the class of CRRA utility functions, is achievable only with the logarithmic utility. \\
The production function is the concave function in quality-adjusted quantity 
\[F(n,z) = [n(z + \alpha_z(1-z))]^\alpha\]
I set the curvature parameter $\alpha=0.85$, consistent with the recent literature on firm dynamics (\textcite{schaal2017}, \textcite{bilal2022}). \\
I work with the CES contact rate function following \textcite{menzio2011}:
\[ p(\theta) = \theta ( 1 + \theta^\gamma)^{-1/\gamma}, \:\:\:\:\:\:\:\:\: q(\theta) = (1+\theta^\gamma)^{-1/\gamma} \]
The idiosyncratic firm-level productivity follows the uniform Markov process as in \textcite{balke2022}, where, with some probability $\lambda_y$, the productivity is redrawn uniformly, and otherwise it remains unchanged. \\
This leaves me with the following parameters requiring estimation: 
\begin{itemize}
    \item Unemployment production $b$
    \item Productivity: correlation $\lambda_y$ and variance $\sigma_y$
    \item Search function: vacancy posting cost $c$, matching function curvature $\gamma$, OJS  search efficiency $\lambda_{jj}$
    \item Firm dynamics: cost of entry $\kappa_e$ and cost of staying $\kappa_f$
    \item Match heterogeneity: ratio of high quality matches upon hiring $q_0$, relative productivity of low quality matches $\alpha_z$
\end{itemize}
\subsection{Moments of interest} %Note: revise the numbers! Some are even made up!!!!
I select four sets of moments to be matched: transition probabilities, wage growth, productivity dynamics, and firm dynamics. \\
First, I document annual transition rates for job-to-job transitions (E2E) as well as the hiring rate U2E. For the hiring rate, I use the proportion of newly hired workers across all observations. The U2E and E2E rates are connected directly to the vacancy posting cost $c$ and the relative OJS search efficiency $\lambda_{jj}$, respectively. I obtain the job-to-job transition rate of $6.3\%$ and the $12.8\%$ proportion of new hires in the economy. \\
I track the distribution of E2U rates across firms based on their productivity. I use this discipline my match heterogeneity parameters as, the lower is the relative productivity of high quality matches $\alpha_z$, the more inclined even highly productive firms will be to layoff workers. Similarly, after a high enough productivity shock firms are incentivized to hire workers. The smaller the rate of high quality new matches $q_0$, the more the firm will potentially choose to fire. \\
%Need the numbers here!!!
I use the wage growth over the first 5 years tenure to discipline the unemployment production $b$, as in \textcite{souchier2022}. The intuition is that, since wages will increase until the marginal profit from the cohort reaches zero, the lower the starting point of the cohort (which depends directly on the outside option of unemployment), the larger will be the initial wage growth. I document the wage growth of $6.1\%$ after 5 years of tenure. \\ %Note: the number's made-up!
To measure overall firm-level productivity dynamics, I look at the standard deviation and persistence of firm-level productivity (excluding the effects of aggregate or sectoral shocks). These moments are meant to discipline the two core productivity parameters of my model: correlation $\lambda_y$ and variance $\sigma_y$. Firm-level productivity is moderately persistent, with the persistence coefficient of 0.79. The standard deviation of fiirm-level productivity is 0.39, the largest of all productivity levels. These estimates are consistent with the comparable estimates of 0.81 and 0.30 provided by \textcite{souchier2022}.\\
Lastly, I document 2 firm-specific statistics. To discipline the entry cost $\kappa_e$, I look target an average fraction of jobs created by opening firms of $6\%$. Finally, because the operating cost $\kappa_f$ affects the exit of firms and, thus, firm selection, I target an average firm size of $32.5$. For comparison, I also documented the average establishment size of $17.9$, which is very similar to $15.6$ from the 2002 US Economic Census. 

\subsection{Untargeted Moments}
As validation of the model, I assess whether it matches the heterogeneity in layoffs and wage passthrough across firms and tenure as documented in Section~\ref{sec:empirics}. \\
Work in progress.
